\begin{center}\begin{tikzpicture}[font=\small,>=stealth,line join=round]
\draw[gray!40,rounded corners=3pt] (0,0) rectangle (14.2,8.3);\node at (7.1,7.700000000000001) {例7：一个在1下方，一个在1上方，无需算小数幂};
\begin{scope}[shift={(1.9571428571428573,2.3576923076923073)},x=2.142857142857143cm,y=1.769230769230769cm]\draw[->,gray] (-0.4,0)--(2.4,0) node[right]{$x$};\draw[->,gray] (0,-0.4)--(0,2.2) node[above]{$y$};\draw[blue,thick,solid,domain=0.05:2,samples=180] plot(\x,{(\x)^(.7)});\draw[violet,thick,dashed,domain=0.05:2,samples=180] plot(\x,{(\x)^(.3)});\fill[orange!85!black] (0.6,0.6993681904144294) circle(2.2pt);\fill[orange!85!black] (1.2,1.0562199684392581) circle(2.2pt);\draw[gray,dashed] (0,1)--(2.1,1);\draw[gray,dotted] (.6,0)--(.6,.6^.7);\draw[gray,dotted] (1.2,0)--(1.2,1.2^.3);\node[left] at (0,1){$1$};\node[below=4pt] at (.6,0){$0.6$};\node[below=4pt] at (1.2,0){$1.2$};\end{scope}\node[anchor=west,align=left,text width=5.0cm] at (8.4,4.0) {蓝：$x^{0.7}$\\[9pt]紫：$x^{0.3}$\\[9pt]$0.6^{0.7}<1$\\[9pt]$1<1.2^{0.3}$\\[9pt]所以$0.6^{0.7}<1.2^{0.3}$};
\node[text width=12.6cm,align=center] at (7.1,.65) {灰色水平虚线为y=1；两个橙点在不同曲线上，分别位于该高度的下方与上方。};
\end{tikzpicture}\end{center}
